\begin{helper}
Let \(p<s\), let \(W\) be a type, and let
\[
  x,y:W\to\mathbb F_2^p
\]
be coordinate maps satisfying \(\langle x(w),y(w)\rangle=1\) for every \(w\in W\).
Define a digraph \(D\) on \(W\) by
\[
  D(a,b)\quad\Longleftrightarrow\quad \langle x(a),y(b)\rangle=0.
\]
Then \(D\) is \(T_s\)-free.
\end{helper}
\begin{proof}
By \(\noderef{TransitiveTournamentFree}\), it is enough to rule out an ordered
\(s\)-tuple \(v_1,\ldots,v_s\) whose earlier vertices all send arcs to later
vertices. Suppose such a tuple exists. Write
\[
  a_i=x(v_i),\qquad b_i=y(v_i)\qquad(1\le i\le s).
\]
The diagonal hypothesis gives
\[
  \langle a_i,b_i\rangle=1
\]
for every \(i\). The arc condition and the definition of \(D\) give
\[
  \langle a_i,b_j\rangle=0
\]
whenever \(i<j\).

We show that \(b_1,\ldots,b_s\) are linearly independent over \(\mathbb F_2\).
Let \(\sum_{j=1}^s c_j b_j=0\). We prove \(c_i=0\) by induction on \(i\).
Assume \(c_j=0\) for all \(j<i\). Taking the dot product with \(a_i\) gives
\[
  0=\left\langle a_i,\sum_{j=1}^s c_j b_j\right\rangle
   =\sum_{j=1}^s c_j\langle a_i,b_j\rangle.
\]
The terms with \(j<i\) vanish by the induction hypothesis, and the terms with
\(j>i\) vanish because \(\langle a_i,b_j\rangle=0\). The remaining term is
\[
  c_i\langle a_i,b_i\rangle=c_i.
\]
Thus \(c_i=0\). Hence all coefficients are zero, proving linear independence.

Therefore \(\mathbb F_2^p\) contains \(s\) linearly independent vectors, so
\[
  s\le \dim_{\mathbb F_2}\mathbb F_2^p=p.
\]
This contradicts \(p<s\). Hence no transitive tournament of size \(s\) exists
in \(D\), so \(D\) is \(T_s\)-free.
\end{proof}
